% Draft LaTeX for the Results section.
% Suggested preamble packages:
% \usepackage{booktabs}
% \usepackage{longtable}
% \usepackage{tabularx}
% \usepackage{pifont}
% \usepackage{xcolor}
% \newcommand{\pass}{\textcolor{green!50!black}{\ding{51}}}
% \newcommand{\fail}[1]{\textcolor{red!65!black}{\ding{55}_{\mathrm{#1}}}}
% \newcommand{\blocked}{\textcolor{gray}{--}}

\subsection{Local model performance}

The local-agent workflow was evaluated on a suite of ODE parameter-estimation
tasks using three locally served Ollama models.  Each trial used the same
user-facing workflow,
\[
\texttt{pfit new} \rightarrow \texttt{pfit check} \rightarrow
\texttt{pfit jax} \rightarrow \texttt{pfit run} \rightarrow
\texttt{pfit diagnose},
\]
with deterministic validation after each LLM-produced artifact.  A run was
counted as successful only if all stages completed, the expected data files and
experiment counts were present, the numerical outputs were finite, the fitting
stage completed, and the diagnosis stage accepted the result.  The comparison
therefore measures end-to-end workflow completion, not merely whether the model
could produce syntactically valid text.

The three-model comparison was run on an NVIDIA H100 80GB HBM3 GPU.  The
Ollama request configuration used a 32k token context window, temperature
0.1, and a 12k token generation budget.  The quantized model artifacts recorded
by the benchmark were approximately 9.0 GB for the 14B model, 17.7 GB for the
Qwen3.8-27B thinking model, and 19.9 GB for the 32B model.  In practice, users
should budget additional memory for the KV cache, Python/JAX runtime, and any
simultaneous numerical fitting workload.  The evaluation therefore used an
80GB accelerator to remove memory pressure from the comparison; smaller cards
can run subsets of the workflow, but large-context local generation is much
more robust on GPUs with at least 32--48GB of memory, and the most comfortable
configuration is an A100/H100-class 80GB device.

Table~\ref{tab:model-completion} summarizes the observed completion behavior.
Two rows were blocked before model evaluation because no complete scientific
problem specification was supplied.  Excluding those blocked rows, Qwen3.8-27B
completed 19 of 21 eligible tasks (90.5\%), compared with 15 of 21 for the
32B baseline (71.4\%) and 13 of 21 for the 14B baseline (61.9\%).  The
Qwen3.8 model was slower in wall-clock time, partly because it is a thinking
model and its reasoning tokens are included in the generation budget, but it
was substantially more reliable on the larger biochemical and multi-experiment
cases.

\begin{longtable}{@{}p{0.27\linewidth}ccc p{0.30\linewidth}@{}}
\caption{End-to-end local workflow completion by model for the retained
user-facing examples.  A check mark denotes completion of all five pfit stages.
A cross indicates the failed stage: N = \texttt{pfit new}, C =
\texttt{pfit check}, and J = \texttt{pfit jax}.  The completion-rate summary in
Table~\ref{tab:model-completion-rates} reports the full benchmark, including
cases that are not retained as examples in \texttt{sessions/}.}
\label{tab:model-completion}\\
\toprule
Case & 32B & 14B & Qwen3.8-27B & Source or model family \\
\midrule
\endfirsthead
\toprule
Case & 32B & 14B & Qwen3.8-27B & Source or model family \\
\midrule
\endhead
ARC fitting & \pass & \pass & \pass & Accelerating-rate calorimetry battery example \\
Boehm STAT5 & \pass & \pass & \pass & STAT5 signaling model \cite{boehm2014stat5} \\
Lotka--Volterra & \pass & \pass & \pass & Predator--prey equations \cite{lotka1925elements,volterra1926fluctuations} \\
MAPK cascade & \pass & \pass & \pass & Kholodenko MAPK cascade \cite{kholodenko2000negative} \\
Oregonator & \pass & \pass & \pass & Field--Noyes/Tyson BZ reaction model \cite{field1974oscillations,tyson1985singular} \\
Piezo Bouc--Wen & \pass & \fail{N} & \pass & Bouc--Wen hysteresis family \cite{bouc1971modele,wen1976method} \\
Robertson & \pass & \pass & \pass & Stiff chemical kinetics model \cite{robertson1966solution} \\
Sliding basepoint & \fail{J} & \pass & \fail{J} & Frictional sliding-basepoint example \\
Theophylline & \pass & \pass & \pass & One-compartment pharmacokinetic example \cite{pinheiro2000mixed} \\
Van der Pol & \pass & \pass & \pass & Nonlinear oscillator \cite{vanderpol1926relaxation} \\
Cascaded tanks & \pass & \fail{N} & \pass & Two-tank process-control example \\
Multi-exponential decay & \pass & \pass & \pass & Synthetic multi-experiment decay example \\
Sneyd IP$_3$ receptor & \fail{N} & \pass & \pass & IP$_3$ receptor model \cite{sneyd2002dynamic} \\
Beer indigoidine & \pass & \pass & \pass & Indigoidine production dataset \cite{beer2014model} \\
Raia IL-13 & \pass & \fail{N} & \pass & IL-13/JAK2/STAT5 lymphoma model \cite{raia2011dynamic} \\
\bottomrule
\end{longtable}

\begin{table}[t]
\centering
\caption{Completion rates for eligible cases.  Blocked rows are excluded
because no complete problem specification was available to any model.}
\label{tab:model-completion-rates}
\begin{tabular}{@{}lrrrr@{}}
\toprule
Model & Pass & Fail & Blocked & Pass / eligible \\
\midrule
Qwen2.5-Coder 32B & 15 & 6 & 2 & 71.4\% \\
Qwen2.5-Coder 14B & 13 & 8 & 2 & 61.9\% \\
Qwen3.8-27B & 19 & 2 & 2 & 90.5\% \\
\bottomrule
\end{tabular}
\end{table}

The main practical conclusion is that local execution is feasible, but only
when the model is strong enough to perform the semantic setup and
Python-to-JAX translation steps reliably.  The 14B model was fast and often
adequate for small canonical systems, but it failed on several larger
multi-experiment or biochemical prompts during \texttt{pfit new}.  The 32B
baseline improved coverage but still failed on several model-construction
tasks.  The Qwen3.8-27B thinking model was slower, yet it was the only model
that completed the Schwen insulin, Armistead sphingolipid, and Fujita EGF
cases in this comparison.  For users who want the workflow to behave like the
remote-agent version on realistic systems biology examples, Qwen3.8-27B is the
preferred local model among those tested.

The completed examples are included in the repository under \texttt{sessions/}.
Each curated session contains the user prompt and data under \texttt{inputs/},
the generated model artifacts under \texttt{generated/}, and at least one
completed fitting run under \texttt{outputs/}.  The \texttt{outputs/}
directory for each retained case also contains one or more
\texttt{fit\_exp*.png} figures showing the fitted trajectory against the data.
These sessions are intended to serve both as examples for users and as
reproducibility evidence for the local-agent workflow.
